% =====================================================================
%  2bat-ud2-3.tex  —  SESSIÓ 3: Quantes combinacions són possibles? (descoberta)
% =====================================================================
\horatitol{Sessió 3 — Quantes combinacions són possibles?}%
{Un cas real: trobar totes les combinacions d'activitats possibles amb un pressupost i unes hores limitats, abans de cap fórmula.}

\subsection{Idea clau}

Avui treballem en grup, a les pissarres, com a gestors d'una entitat. Tenim un
\textbf{pressupost màxim} i un \textbf{nombre màxim d'hores} de monitoratge, i hem
d'organitzar \textbf{dos} tipus d'activitat. El repte: trobar quines combinacions són
possibles sense passar-nos de res. No us dono cap fórmula: l'heu de descobrir.

\begin{idea}
\textbf{El cas.} Una entitat vol organitzar \textbf{tallers} ($x$) i \textbf{sortides}
($y$). Cada activitat consumeix diners i hores de monitoratge:
\begin{center}
\begin{tabular}{lcc}
\toprule
 & Cost (\euro) & Hores de monitoratge \\
\midrule
Cada taller ($x$) & $100$ & $2$ \\
Cada sortida ($y$) & $200$ & $5$ \\
\bottomrule
\end{tabular}
\end{center}
Disposem com a màxim de \textbf{$\num{1200}\eur$} i de \textbf{$30$ hores}.
\textbf{La pregunta del dia:} quines combinacions $(x,y)$ són possibles? N'hi ha una
o moltes? Si les dibuixéssiu totes, quina forma tindrien?
\end{idea}

\subsection{Exemples}

\begin{exemple}[provar una combinació]
Provem $3$ tallers i $2$ sortides, és a dir $(x,y)=(3,2)$:
\begin{itemize}[leftmargin=1.4em, itemsep=2pt]
  \item \textbf{Cost:} $100\cdot 3+200\cdot 2=300+400=700\eur\le\num{1200}$: d'acord.
  \item \textbf{Hores:} $2\cdot 3+5\cdot 2=6+10=16\le 30$: d'acord.
\end{itemize}
La combinació $(3,2)$ \textbf{és possible}: compleix totes dues condicions.
\end{exemple}

\begin{exemple}[una combinació que no es pot]
Provem $(8,4)$: $8$ tallers i $4$ sortides.
\begin{itemize}[leftmargin=1.4em, itemsep=2pt]
  \item \textbf{Cost:} $100\cdot 8+200\cdot 4=800+800=\num{1600}\eur$, que
        \textbf{supera} els $\num{1200}\eur$. Ja no cal mirar les hores: \textbf{no} es pot.
\end{itemize}
N'hi ha prou que falli \emph{una} condició perquè la combinació quedi descartada.
\end{exemple}

\subsection{Exercicis resolts}

\begin{exresolt}[traduir les condicions a inequacions]
Escriu les condicions del cas (pressupost i hores) com a inequacions amb dues
variables, sense oblidar que no es poden fer activitats negatives.
\solucio
\textbf{Pressupost:} cada taller costa $100\eur$ i cada sortida $200\eur$, i el total
no pot superar $\num{1200}$:
\[
100x+200y\le \num{1200}.
\]
\textbf{Hores:} cada taller són $2$ h i cada sortida $5$ h, com a màxim $30$:
\[
2x+5y\le 30.
\]
\textbf{No-negativitat:} $x\ge 0$ i $y\ge 0$ (no hi ha activitats negatives).
\resposta{$100x+200y\le\num{1200}$, \;$2x+5y\le 30$, \;$x\ge 0$, \;$y\ge 0$.}
\end{exresolt}

\begin{exresolt}[comprovar una combinació amb totes les condicions]
És possible la combinació $(6,3)$ ($6$ tallers i $3$ sortides)?
\solucio
Comprovem les dues condicions:
\begin{itemize}[leftmargin=1.4em, itemsep=1pt]
  \item Cost: $100\cdot 6+200\cdot 3=600+600=\num{1200}\le\num{1200}$: d'acord (just al límit).
  \item Hores: $2\cdot 6+5\cdot 3=12+15=27\le 30$: d'acord.
\end{itemize}
Compleix totes dues (i $x,y\ge 0$): la combinació $(6,3)$ \textbf{és possible}, just al
límit del pressupost.
\resposta{Sí: cost $\num{1200}$ (just al límit) i $27$ hores; compleix tot.}
\end{exresolt}

\subsection{Exercicis per fer}

\noindent\textit{Treballeu en grup. Useu el cas: cada taller $100\eur$ i $2$ h; cada
sortida $200\eur$ i $5$ h; màxims $\num{1200}\eur$ i $30$ h.}

\begin{exercicis}
\begin{exlist}
  \item Comprova si cada combinació $(x,y)$ és possible (mira cost \emph{i} hores):
        \begin{enumerate}[label=\alph*), leftmargin=1.6em, topsep=2pt]
          \item $(4,2)$ \hfill \item $(2,5)$ \hfill \item $(10,1)$
        \end{enumerate}
  \item Escriu les \textbf{quatre} condicions del problema com a inequacions (les dues
        de recursos i les dues de no-negativitat).
  \item Troba \textbf{tres} combinacions diferents que siguin possibles i \textbf{una}
        que no ho sigui, justificant cada cas amb els càlculs.
  \item Si l'entitat només fes \textbf{sortides} ($x=0$), quantes en podria fer com a
        màxim? (Mira quina condició és la més restrictiva: pressupost o hores.)
  \item Si només fes \textbf{tallers} ($y=0$), quants en podria fer com a màxim? Quina
        condició mana en aquest cas?
  \item \textbf{(Descoberta)} A la pissarra, marqueu uns quants punts $(x,y)$
        possibles i uns quants d'impossibles. Quina \textbf{forma} sembla tenir la
        zona de les combinacions possibles? Per què creieu que les solucions no són
        casos solts sinó tota una regió?
\end{exlist}
\end{exercicis}

% ---------------------------------------------------------------------
\fullsolucions{Sessió 3}
\begin{solucions}
\begin{sollist}
  \item (a) cost $400+400=800$, hores $8+10=18$: \textbf{possible}. \quad
        (b) cost $200+\num{1000}=\num{1200}$, hores $4+25=29$: \textbf{possible}
        (just al límit de cost). \quad
        (c) cost $\num{1000}+200=\num{1200}$, hores $20+5=25$: \textbf{possible}.
  \item $100x+200y\le\num{1200}$; \;$2x+5y\le 30$; \;$x\ge 0$; \;$y\ge 0$.
  \item Resposta oberta. Exemples possibles: $(3,2)$, $(6,3)$, $(0,6)$; impossible:
        $(8,4)$ (cost $\num{1600}>\num{1200}$).
  \item Amb $x=0$: pressupost $200y\le\num{1200}\Rightarrow y\le 6$; hores
        $5y\le 30\Rightarrow y\le 6$. Totes dues donen $6$: com a màxim \textbf{$6$
        sortides}.
  \item Amb $y=0$: pressupost $100x\le\num{1200}\Rightarrow x\le 12$; hores
        $2x\le 30\Rightarrow x\le 15$. Mana la més petita: com a màxim \textbf{$12$
        tallers} (el pressupost és la condició restrictiva).
  \item Resposta oberta. La zona possible té forma de \textbf{polígon} (regió del
        pla). No són casos solts perquè qualsevol punt dins de la zona compleix totes
        les condicions: hi ha infinites combinacions vàlides.
\end{sollist}
\end{solucions}
